\section{实验}

\subsection{实验设置}
\paragraph{数据集与划分。}
训练阶段使用 Ref-VOS 训练集进行监督学习；评测阶段在验证集上报告分割指标。
每个表达式对应一组候选实例，模型在组内完成排序判别。

\paragraph{实现细节。}
基座模型采用 Qwen3-VL，并通过 LoRA 进行参数高效微调。
视觉编码器冻结，LoRA 注入到语言侧注意力与前馈层关键线性投影。
训练采用 listwise 样本组织（每个 batch 为一个表达式组），优化器为 AdamW，支持混合精度与梯度检查点。
候选掩码监督由离线阶段预计算并缓存。

\paragraph{输入与打分。}
每个候选由多帧 tube 表示；将候选掩码叠加到原图后，与统一 Yes/No 提示词共同输入模型。
模型输出 $P(\text{Yes})$ 作为候选匹配分数，并在同组候选内排序。

\paragraph{监督与损失。}
离线阶段预计算候选与 GT 的 IoU，依据阈值 $\tau_{pos},\tau_{neg}$ 构建正负候选集合。
训练目标由 listwise 排序损失与软掩码对齐损失组成。

\paragraph{评测指标。}
采用区域相似度 $\mathcal{J}$（IoU）、边界准确度 $\mathcal{F}$，以及二者均值 $\mathcal{J}\&\mathcal{F}$。
默认以 expression-level 平均结果作为主指标，并补充 video-level 统计。

\subsection{主结果}
\begin{table}[t]
\centering
\caption{验证集主结果对比。}
\begin{tabular}{lccc}
\hline
方法 & $\mathcal{J}$ & $\mathcal{F}$ & $\mathcal{J}\&\mathcal{F}$ \\
\hline
原始 Qwen3-VL + CRF-SAM3 & -- & -- & -- \\
本文方法（Listwise + Soft Alignment + LoRA） & -- & -- & -- \\
\hline
\end{tabular}
\label{tab:main_cn}
\end{table}

\subsection{消融实验}
\begin{table}[t]
\centering
\caption{关键模块消融。}
\begin{tabular}{lccc}
\hline
配置 & $\mathcal{J}$ & $\mathcal{F}$ & $\mathcal{J}\&\mathcal{F}$ \\
\hline
Pointwise 监督 & -- & -- & -- \\
+ Listwise 排序 & -- & -- & -- \\
+ 软掩码对齐约束 & -- & -- & -- \\
+ LoRA 微调（完整设置） & -- & -- & -- \\
\hline
\end{tabular}
\label{tab:ablation_cn}
\end{table}

\subsection{结果分析}
\paragraph{Listwise 作用。}
Listwise 损失直接优化同组候选排序，能够缓解 pointwise 训练与推理排序目标不一致问题。

\paragraph{软掩码对齐作用。}
软对齐损失提升候选掩码区域内证据质量占比，抑制背景特征捷径，从而改善实例级判别。

\paragraph{效率与稳定性。}
离线缓存将候选召回与 VLM 微调解耦，显著降低训练期显存峰值并提升迭代稳定性。

\subsection{可复现性说明}
\begin{itemize}
\item 训练与推理须保持提示词模板及打分定义一致（均采用 Yes/No 概率打分）。
\item 评测时需保证每个表达式覆盖完整视频帧，缺失帧按背景预测处理。
\item 进行预训练与微调模型对比时，候选生成与排序配置需严格一致。
\end{itemize}
